\subsection{Complex locally convex spaces}

To say that a subset A of a complex vector space E is balanced means that, for all \(x \in A\) , we have \(\rho x \in A\) for \(0 \leqslant \rho \leqslant 1\) and \(e^{i\vartheta}x \in A\) for all real \(\vartheta\).

We say that a set A of E is convex if it is convex in the real space \(E_{0}\) which underlies E. In order that a convex set \(A \neq \varnothing\) of E be balanced, it is sufficient that \(e^{i\vartheta}A \subset A\) for all real \(\vartheta\) ; for this implies firstly that -A = A; as A is convex, we see that 0 belongs to A and thus \(\rho A \subset A\) for \(0 \leqslant \rho \leqslant 1\) .

Let E be a complex topological vector space. The smallest balanced convex (resp. closed balanced convex) set containing a set A of E is called the balanced convex envelope (resp. balanced closed convex envelope) of A; the balanced closed convex envelope of A is the closure of the balanced convex envelope of A. This last is the convex envelope of the union of the sets \(e^{i\theta}A\) ; we can therefore define it as the set of linear sums \(\sum_{i}\lambda_{i}x_{i}\) , when \((x_{i})\) is any finite family of points of A, and \((\lambda_{i})\) a family of complex numbers such that \(\sum_{i}|\lambda_{i}| \leqslant 1\) . If A is precompact so also is its balanced envelope (I, p.~6, prop. 3).

We say that a complex topological vector space E is locally convex if the real underlying topological vector space \(E_{0}\) is locally convex, that is to say if every neighbourhood of 0 in E contains a convex neighbourhood of 0; a topology T on E is locally convex if it is compatible with the vector space structure of E (relative to C) and if E, with topology T, is locally convex. As in this case every closed convex neighbourhood V of 0 contains a balanced neighbourhood W of (I, p.~7, prop. 4), we see that V also contains U, the balanced closed convex envelope of W; in other words the balanced, closed, convex neighbourhoods of 0 form a fundamental system of neighbourhoods of 0 in E, invariant under every homothety of ratio \(\neq 0\) .

Conversely, let E be a complex vector space and let S be a filter base on E formed by absorbent, balanced convex sets. We know then (II, p.~23, prop. 1) that the set B, of the transforms of the sets of S by homotheties of ratio \textgreater{} 0, is a fundamental system of neighbourhoods of 0 for a locally convex topology T on the real vector space \(E_{0}\) underlying E. Further, as the sets of B are balanced, they are invariant under every homothety \(x \mapsto e^{i\vartheta}x\) , which shows that T is compatible with the vector space structure of E (over C) (I, p.~7, prop. 4).

Every locally convex topology on a complex vector space E can be defined by a set of semi-norms, for the gauge of an open balanced convex neighbourhood of 0 is a semi-norm on E.

The ideas and results for real locally convex spaces detailed in II, p.~25 to 36, extend to complex locally convex spaces with no modification other than the replacement of symmetric convex sets by balanced convex sets.

A complex locally convex space is a Fréchet space if it is metrisable and complete.

